\section{问题重述与分析}
\subsection{任务及决策层次}
研究对象为含光伏和储能设备的微网。微网不向外部电网售电，日常用电由光伏、已经约定的外网购电及储能共同承担；电力不足时，可以按对应电价的五倍紧急购电。储能额定容量为12000\,kWh，允许储电量为1200--10800\,kWh，充放电最大功率均为5000\,kW。

题目要求解决四类逐步增加信息与决策自由度的问题：其一，在预测曲线确定且当天储电量闭合的条件下，求一天的最低购电费用；其二，仅以历史运行数据制定每日零时计划，并在真实负荷与光伏实现时调整储能；其三，结合每日零时、6时、12时和18时发布的光伏预报，允许对尚未交付的购电计划进行调整，分析新预报的价值；其四，将上述策略推广至波动电价。除全年结果外，还需给出指定时段与四个指定日期的详细调度表。

问题一是确定性资源配置问题，可以直接建立线性规划。问题二至四同时涉及预测和控制：即使某次规划在预测输入下最优，也不能据此推出有预测误差时的全年策略最优。因此本文将\emph{预测、常规购电规划和实时执行}分为三个环节，分别约束信息的可用时刻、物理动作和费用核算。储能跨时间耦合又使预测误差的经济影响依赖库存及电价，故预测精度只能作为辅助指标，主要评价量应为真实结算费用。

\subsection{数据整理与时间约定}
附件一包含144个时点的电价及典型负荷、光伏预测，附件二和附件四包含2025年365天、每天144个时点的实际功率和电价，附件三包含每日4次、每次未来24小时的光伏预报。日期连续，所需数值无缺失、负数或非有限值。输出评价区间为2025年2月1日至12月31日，共334天、48096个十分钟区间；1月用于历史初始化、候选参数验证和控制热启动。

输入记录00:10被解释为00:00--00:10区间的平均代表功率，因而单段电量等于该行功率乘以$\Delta=1/6$小时。这一解释与输入的右端时标一致；由于题目未区分瞬时采样与平均功率，另对问题一作梯形积分敏感性检验。原结果模板的标签整体向后偏移十分钟，故输出统一标为00:00--00:10至23:50--24:00，保持数据次序不变。

附件三第$k$个预报值对应\emph{发布时间之后第$k$小时}的目标，不能按当日第$k$小时读取。本文使用目标整点间线性插值；发布时间的插值起点取当时已经完成区间的光伏观测。新发布的6时预报最早影响6:00--6:10，不能回写已经交付的区间。

附件一的典型负荷与附件二全年同刻均值几乎一致。如果把这一曲线作为问题二至四的预测先验，就会隐含使用未来统计量。因此后续问题仅引用附件一已给定的固定电价，不使用其典型负荷或光伏作为历史先验。问题一则依题意直接使用附件一的给定预测。

\subsection{方法背景与本文定位}
储能调控需要在不同电价和预测误差之间分配有限库存。分时电价下光伏储能经济调度及日前、日内协调优化已有相关研究\cite{wu2018tou,xiao2019mpc}。模型预测控制在每次获得新观测后，重新预测并优化未来时域，仅实施首个动作；确定性滚动控制与多情景鲁棒控制具有不同的不确定性处理方式\cite{wytock2017dem}。动态电价研究也强调预测驱动的可实施策略应与全知性能参考区分\cite{perez2023hem}。

本文采用容易核验的风险分位规划与显式储能反馈，而不把经验分位包装为联合概率约束、条件风险价值或分布鲁棒最优。通过历史预测基准、严格预报信息对照、替代交易合约以及价格感知反馈比较，检验结论对建模选择的依赖性。模型评价的重点是费用、信息边界和状态连续性，而不是算法名称的复杂程度。
